\documentclass[11pt]{article}

\usepackage[margin=1in]{geometry}
\usepackage{parskip}
\usepackage{listings}
\usepackage{xcolor}
\usepackage{booktabs}
\usepackage{longtable}
\usepackage{hyperref}
\hypersetup{colorlinks=true,linkcolor=blue,urlcolor=blue}

\newcommand{\code}[1]{\texttt{\detokenize{#1}}}

\lstset{
  basicstyle=\ttfamily\footnotesize,
  breaklines=true,
  columns=fullflexible,
  frame=single,
  framesep=4pt,
  backgroundcolor=\color{gray!8},
  showstringspaces=false,
}

\title{Is There an API to Obtain \emph{All} of a Swimmer's Previous
Times from the USA~Swimming Database?}
\author{Findings from the \code{times-api} OpenAPI specification}
\date{\today}

\begin{document}
\maketitle

\section{Short answer}

\textbf{Yes.} The USA~Swimming data-hub back end (\code{times-api},
serving \url{https://data.usaswimming.org/datahub/}) exposes a
first-party REST endpoint whose explicit purpose is to return
\emph{every} recorded swim for a member, not merely lifetime bests:

\begin{quote}
\code{POST /swims/TimesSearch/GetAllTimesForFilters}
\end{quote}

It is complemented by several other endpoints that return full history
sliced by meet, by season, or in a single flat per-person list. This
report catalogues each of them, gives request and response shapes, and
notes how they differ from the best-times endpoints the existing
\code{swim-times} program currently uses.

\section{Evidence and method}

The findings are drawn from the machine-readable OpenAPI~3.0.4
specification shipped with this project as
\code{swim2/swagger.json} (399\,KB, \code{title: times-api},
\code{version: 1.0}). This is the authoritative contract published by
the service itself. The document defines 120+ operations under the
\code{/swims/} prefix; the ones relevant to retrieving an individual
swimmer's historical times are extracted below.

The spec declares \emph{no} \code{securitySchemes} and \emph{no} global
\code{security} requirement, which is consistent with the project's
existing experience that the times endpoints are reachable by anonymous
callers who supply the three data-hub headers (\code{AppName: DataHub},
\code{Usas-Sub-Id: Anonymous}, and a client-minted \code{Device-Id})
rather than a bearer token.

Note: the response shapes below are transcribed from the specification's
schema definitions. A live call to confirm field-level behaviour was not
executed for this document; where empirical confirmation matters (e.g.\
exact pagination limits) it should be verified against the live service.

\section{The primary answer: \code{GetAllTimesForFilters}}

\begin{center}
\begin{tabular}{@{}ll@{}}
\toprule
Method \& path & \code{POST /swims/TimesSearch/GetAllTimesForFilters} \\
Request body   & \code{SFTimesSearchFilters} (JSON) \\
Response       & array of \code{SFTimesSearchResults} \\
\bottomrule
\end{tabular}
\end{center}

Unlike the best-times calls, this endpoint is a general times \emph{search}:
its filter object lets the caller pin a single \code{memberId} and then
optionally narrow (or not) by event, course, season, date range, age
range, and time-standard type. Leaving the narrowing fields null returns
the member's complete history for the requested events/course.

\textbf{Request --- \code{SFTimesSearchFilters}:}
\begin{lstlisting}
{
  "memberId":         string,   // the swimmer
  "events":           string,   // e.g. "50 FR"; null/broad => all
  "course":           string,   // "SCY" | "LCM" | "SCM"
  "seasonKey":        integer,  // null => all seasons
  "startDate":        string,   // null => no lower bound
  "endDate":          string,   // null => no upper bound
  "minAge":           integer,
  "maxAge":           integer,
  "timeStandardType": string,
  "sortBy":           string
}
\end{lstlisting}

\textbf{Response --- array of \code{SFTimesSearchResults}, one element
per swim:}
\begin{lstlisting}
{
  "swimTimeId":   integer,
  "meetId":       integer,
  "meetName":     string,
  "eventCode":    string,
  "courseCode":   string,
  "swimDate":     string,
  "swimmerAge":   integer,
  "powerPoints":  integer,
  "swimTime":     string,
  "timeStandard": string,
  "clubName":     string
}
\end{lstlisting}

Because each element carries its own \code{swimDate}, \code{meetName},
and \code{timeStandard}, and the array is not collapsed to one row per
event, this is the endpoint that yields a swimmer's \emph{full
historical progression} (every heat/final on record), which is exactly
what the question asks for.

\section{Other endpoints that return full history}

The service offers several complementary ways to obtain all of a
swimmer's times, differing in how the results are sliced.

\subsection{\code{GET /swims/TimesSearch/MemberTimes}}

A query-string variant with an explicit \code{bestTimesOnly} toggle ---
setting it false returns all times.

\begin{lstlisting}
GET /swims/TimesSearch/MemberTimes
    ?personId={id}
    &organization=...
    &eventCourse=SCY|LCM|SCM
    &session=...
    &season=...
    &bestTimesOnly=false        // false => full history
-> array of SFMemberTimes
   { fullName, eventCode, swimTime, swimTimeAdj, swimmerAge,
     powerPoints, timeStandardType, meetName, lscClub, swimDate }
\end{lstlisting}

The presence of a \code{bestTimesOnly} flag is direct evidence that the
same backing query can return either the reduced best-times view or the
complete set.

\subsection{Meet-by-meet drill down}

Two endpoints together reconstruct a swimmer's whole career meet by meet:

\begin{lstlisting}
GET /swims/TimesSearch/GetSwimmerMeets/{memberId}
-> array of SFMeet
   { meetId, meetName, meetType, seasonYear, season, minDate, maxDate,
     swims, lifetimeBestTimes, seasonBestTimes, courseCode, meetDates }

GET /swims/TimesSearch/GetSwimmerMeetTimes/{memberId}/{meetId}
-> array of SFSwimmerMeetTime
   { swimTimeId, eventCode, sessionName, swimTime, timeStandard,
     timeDrop, finishPosition }
\end{lstlisting}

\code{GetSwimmerMeets} lists every meet the swimmer has competed in
(with a \code{swims} count per meet); iterating \code{GetSwimmerMeetTimes}
over those meet ids yields every individual swim. This is the most
faithful reproduction of the ``meet results'' view on the public site,
and uniquely exposes per-swim \code{finishPosition} and \code{timeDrop}.

\subsection{Flat per-person list}

\begin{lstlisting}
GET /swims/Time/Person/{personId}
-> array of PersonSwimTime
   { swimTimeId, swimTime, competitor, eventName, eventCourse, swimDate,
     session, finishPosition, distance, meetName, orgUnitName,
     standardName, splits[] }

GET /swims/Time/Mobile/{personId}
-> array of PersonSwimTimeMobile   // leaner mobile projection
\end{lstlisting}

\code{Time/Person} returns a single flat array of every swim for the
person, and notably includes \code{splits[]} (per-length split times)
for each swim --- the richest per-swim detail of any endpoint here.

\subsection{Individual swim + splits}

\begin{lstlisting}
GET /swims/TimesSearch/AllTimes/SwimTime/{swimTimeId}
GET /swims/TimesSearch/SwimmerMeets/SwimTime/{swimTimeId}
\end{lstlisting}

Given a \code{swimTimeId} harvested from any list above, these fetch a
single swim's full detail, including splits (\code{SFSwimSplit}:
cumulative distance, cumulative split time, split time).

\section{Contrast with the best-times endpoints currently used}

The existing \code{swim-times} program uses only the two best-times
endpoints, which by design collapse history to one row per event:

\begin{center}
\begin{longtable}{@{}p{0.34\textwidth}p{0.60\textwidth}@{}}
\toprule
\textbf{Endpoint} & \textbf{What it returns} \\
\midrule
\endhead
\code{GET .../GetBestTimesForMember/\{id\}} &
  Lifetime best per event; stroke, distance, course, time \emph{only}
  (no date, meet, or standard). \code{SFSwimmerBestTime}. \\
\addlinespace
\code{POST .../TimesSearch/BestTimes} &
  Lifetime best per stroke/distance pair \emph{with} date, meet,
  standard, splits, finish position. \code{SFBestTimesSearchResults}.
  One row per event. \\
\midrule
\code{POST .../GetAllTimesForFilters} &
  \textbf{All} swims matching the filter --- every historical swim, each
  with its own date/meet/standard. \code{SFTimesSearchResults}. \\
\bottomrule
\end{longtable}
\end{center}

The best-times endpoints answer ``what is this swimmer's fastest ever
50~free?''. \code{GetAllTimesForFilters} (and its siblings) answer
``show me every 50~free this swimmer has ever recorded, when, and where''
--- the full progression.

\section{Recommendation}

To obtain all of a swimmer's previous times programmatically, use
\code{POST /swims/TimesSearch/GetAllTimesForFilters} with an
\code{SFTimesSearchFilters} body that sets \code{memberId} (and
optionally \code{course}/\code{events}) while leaving the date, season,
and age filters null. For per-swim split detail, follow up with
\code{GET /swims/Time/Person/\{personId\}} or resolve individual
\code{swimTimeId}s. These calls should be reachable with the same
anonymous data-hub headers (\code{AppName}, \code{Usas-Sub-Id},
\code{Device-Id}) the project already mints for the best-times calls,
since the spec defines no authentication scheme; this should be
confirmed with a live request before relying on it in production.

\end{document}
